\section{物理训练如何塑造 AI 研究品味}

姚顺宇的物理经历不是访谈的花絮，而是理解他研究判断的重要背景。他本科做凝聚态理论，后来做理论高能、量子信息与黑洞相关方向。谈到非厄米系统时，他强调本科时期的研究和现在做 AI 有相似性：先有一个理解或想法，再设计数值实验验证，最后把想法落实成 pipeline。

\begin{figure}[H]
\centering
\includegraphics[width=0.88\textwidth]{frame-physics-nonhermitian.jpg}
\caption{物理经历章节：从凝聚态、非厄米系统到 AI 实验方法。\protect\footnotemark}
\end{figure}
\footnotetext{视频画面时间区间：01:28:25--01:28:55。该帧用于标记物理背景讨论。}

\subsection{非厄米系统的直觉版解释}

在普通量子力学里，哈密顿量通常是厄米的，这保证能量本征值是实数，系统演化保持概率守恒。非厄米系统则常用来描述开放系统、耗散、增益、测量或有效动力学。对 AI 笔记读者来说，不必深入物理细节，但要理解它训练出的研究习惯：面对复杂系统时，先构造可检验模型，再用数值实验逼近理解。

\begin{knowledgebox}{从物理到 AI 的迁移}
物理训练强调：定义问题、建立简化模型、设计实验、观察异常、修正理论。大模型训练也类似：提出行为假设，设计数据和训练 pipeline，用实验观察 loss、benchmark、用户行为和失败案例，再回到假设。差异在于，AI 系统更工程化、更依赖组织协作。
\end{knowledgebox}

\subsection{“挑战不擅长的事”}

访谈中，姚顺宇多次说自己喜欢挑战不擅长的事。离开物理并不是因为物理无趣，而是他觉得自己在那个方向已经看到很多比自己聪明的人，也想去挑战一个新领域。这个选择解释了他后来去 Anthropic、再去 Gemini 的动机：不是追求头衔，而是寻找能迫使自己学习新东西的环境。

\begin{warningbox}{不要把跨界神话化}
从物理转 AI 并不自动意味着“降维打击”。姚顺宇反而多次说 AI 不一定需要最抽象的聪明，而需要靠谱、细致、负责。跨界带来的价值主要是研究品味和实验习惯，而不是天然优越感。
\end{warningbox}

\subsection{本章小结}

物理背景给姚顺宇提供的是一种处理复杂系统的方式：从想法到实验，从实验到理解，再到 pipeline。这个方法迁移到 AI 后，表现为对问题定义、数据构造和实验验证的重视。

